\section{Introduction}

A version control system records changes so they can be revisited and reversed.
Every widely adopted system records lines of code---whether as
snapshots~\cite{git2026} or as patches to those
lines~\cite{roundy2005,pijul2024}---and groups them into commits a developer
names. This agreement dates to the earliest systems and has survived every
redesign since, for a reason that was sound from the start: the developer
authored the code, held at least a rough model of what changed and
why~\cite{naur1985,ko2006}, and could therefore supply the decomposition
whenever a tool asked for it. Developers now specify intent in natural language and models translate it into
edits across files~\cite{sarkar2025}. Three controlled field experiments show
this shift produces 26\% more code per week with no measurable quality
decline~\cite{cui2025}. The condition that made the agreement sound---that the
developer holds a model of what changed and why---no longer holds reliably. A
developer who types three sentences in an hour and receives fourteen functions
across five files cannot supply the decomposition a tool asks for, because they
never held it. We argue that the design choice should be renegotiated.

Three consequences of this shift make the mismatch concrete. First, the unit
that version control addresses, the line, is no longer stable enough to serve as
an anchor. A model that generates a function in one response and regenerates it
with different formatting or variable names in the next has changed every line
while changing nothing a developer would name. Across 211 million changed lines
from public repositories authored between 2020 and 2024, the share associated
with refactoring---lines moved into a better position rather than written
afresh---fell from a quarter to under a tenth, while code revised within two
weeks of being written roughly doubled~\cite{gitclear2025}. A line is now
replaced faster than it settles into a position worth recording, and the
developer wants to revisit a feature as it stood two requests ago---a unit
the line-level record cannot express.

Second, the developer does not hold a detailed model of what the code
says~\cite{sarkar2022}. A natural language prompt specifies an outcome
but not the logic that implements it, so the implementation contains decisions
the model made and the developer has not reviewed. Barke et al.\ observed
programmers reading generated code for shape rather than line by
line~\cite{barke2023}, Tang et al.\ found that developers often do not notice
which code a model wrote and delete it rather than repair it~\cite{tang2024},
and among 410 developers surveyed by Liang et al.\ understanding generated code
was one of the most frequently reported
difficulties~\cite{liang2024}. The consequence for version
control is that the developer who searches a history for a prior state must
describe it in terms they never acquired. This is the vocabulary problem Furnas
et al.\ measured in information retrieval~\cite{furnas1987}, now appearing inside
the repository. Nguyen and Glassman observed it directly: beginning programmers
and code models systematically misread each other's vocabulary, so a developer's
search terms and the model's naming choices diverge even when both refer to the
same behaviour~\cite{nguyen2024}. The developer searches with vocabulary they
never learned and finds whatever happens to match instead of the version that
served a stated purpose~\cite{codoban2015}.

Third, the unit the developer thinks in has moved from a location in the code to
a purpose the code serves. A prompt names what the software should do (rather
than which file to edit), so the developer's mental model of their own work is
organised by feature: the retry policy, the export endpoint, the rate limiter. This is the
level the program comprehension literature has always found people work at:
Biggerstaff et al.\ called recovering which functions implement a named concern
the concept assignment problem~\cite{biggerstaff1993}, and the feature location
field exists because answering it by reading code is
expensive~\cite{dit2013}. The commit is the closest unit a version control system
offers, and it does not correspond to this level, since Herzig and Zeller found
changes serving several unrelated purposes throughout the histories of five open
source Java projects~\cite{herzig2013}. An agent session produces exactly that
structure, because a commit ends where the developer stopped to type a command
and a request ends somewhere else.

The response we describe is a version control system that records at the grain
the developer thinks at. Instead of recording lines inside a commit, we parse
each file into its functions and classes~\cite{treesitter} and record one edit
per changed function, together with what that edit was built on. We group those
functions into features from how they change and how they refer to one
another~\cite{zimmermann2004,ying2004,traag2019}, and label each group with the
words the developer and the agent used while the work was happening. A version
of the code is then a set of these edits, and the files on disk are rebuilt
from whichever set is included, so a feature can be taken out or put back by
naming it rather than by naming the lines it produced. None of the individual techniques is new. Change coupling mining, community
detection, feature location, and set-based version control each have a
literature. What is new is composing them into a single tool that a developer
operates daily, because that composition introduces interactions none of the
original designs anticipated. A function identity tracker that loses a rename
breaks the dependency chain a revert depends on. A grouping inference that absorbs
a shared helper pulls unrelated work into a single feature. A layout record (the
whitespace between functions) that becomes individually addressable acquires
failure modes the entity record does not have, and 68\% of the violations we found
in ten thousand operations entered through that seam. The contribution is in making
the composition survive daily use on a real codebase, in measuring where it fails
honestly, and in reporting both the design's costs and the conditions under which
it should not be adopted.

For example, a service that answers price quotes for a trading dashboard has a
handler, a database layer, and a set of row accessors. On Monday the developer
prompts \cmd{cache the price lookup, it takes 400ms and the price barely moves},
and the agent writes \sym{cache.py::get\_cached} and \sym{cache.py::\_ttl}, gives
\sym{db.py::query} a parameter for how long a cached price stays good so the cache
can pass one through, and changes \sym{api.py::handle} to consult the cache before
the database. That is four edits to four functions in three files, and the two in
\sym{cache.py} share no line with the one in \sym{api.py}, yet all four are filed
under one name. Over the next two days two further prompts change
\sym{api.py::handle} again, one adding a rate limit and one adding a retry policy,
and while it is adding the retry the agent goes back and reworks two of the cache's
own functions, which brings the price cache to six edits. A colleague adds a metric
for quote latency. Three weeks later the developer
finds the dashboard showing a price sixty seconds stale, so they type
\cmd{sgt revert "price cache"}. \sgt{} prints the six edits that would come out,
names \sym{api.py::handle} as the one function where later work was built on
something being removed, shows the text each change would produce, and asks before
touching anything. The rate limit and the retry policy stay, because their edits
are still included; the colleague's metric stays, because it was never in the set
coming out; and the parameter leaves both \sym{db.py::query} and
\sym{api.py::handle}, because both were in it. When the developer decides a week
later that the cache was worth keeping after all, \cmd{sgt restore} returns the
files byte for byte, because putting the six edits back is the same set arithmetic
as taking them out and nothing is replayed as a patch.
Figure~\ref{fig:two-records} draws these three requests arriving in a single
sitting, under both records, and Section~\ref{sec:walkthrough} follows this
repository through five sittings.

Three prices are paid for this, all of them in the first hour of use rather than
in some later edge case. Two edits to the same function made from the same
starting point are recorded as competing versions that a person has to decide
between, including the pair at opposite ends of a fifty line function that touch
no common line and that git would have merged without comment. The grouping of
edits into named requests is inferred, so it is a guess and it will sometimes be
wrong, and correcting it costs the developer a rename, a merge, or a split, none
of which moves a recorded edit. \sgt{} reads Python, TypeScript, and TSX and
versions every other file whole, at the grain git already uses, so a repository
whose interesting logic is written in Go gets nothing from \sgt{} today.

The paper makes four contributions. We characterise the mismatch between the unit
a developer using an agent decides in and the unit a version control system
records, one attempted operation at a time (Section~\ref{sec:scenarios}), in the
style of conceptual analysis Pérez De Rosso and Jackson applied to git
itself~\cite{perezderosso2013,perezderosso2016}. We describe \sgt{}'s design as
four commitments, say what each one buys and what it costs, and run the whole set
of operations on one repository (Sections~\ref{sec:design}
and~\ref{sec:walkthrough}). We state which observation would show each of our
claims to be wrong, and pre-register the study that would produce it
(Section~\ref{sec:study}). We report what fifty-seven days of \sgt{} versioning its
own construction taught us about where the design breaks
(Section~\ref{sec:discussion}), and what three pilots taught us about the distance between a property that holds
across 10{,}237 operations and a property a new user can rely on
(Section~\ref{sec:pilot}).
